% Options for packages loaded elsewhere
\PassOptionsToPackage{unicode}{hyperref}
\PassOptionsToPackage{hyphens}{url}
\PassOptionsToPackage{dvipsnames,svgnames,x11names}{xcolor}
\documentclass[
  american,
  11pt,
  letterpaper,
]{article}
\usepackage{xcolor}
\usepackage[margin=1in]{geometry}
\usepackage{amsmath,amssymb}
\setcounter{secnumdepth}{-\maxdimen} % remove section numbering
\usepackage{iftex}
\ifPDFTeX
  \usepackage[T1]{fontenc}
  \usepackage[utf8]{inputenc}
  \usepackage{textcomp} % provide euro and other symbols
\else % if luatex or xetex
  \usepackage{fontspec}
  \defaultfontfeatures{Scale=MatchLowercase}
  \defaultfontfeatures[\rmfamily]{Ligatures=TeX,Scale=1}
\fi
\usepackage{lmodern}
\ifPDFTeX\else
  % xetex/luatex font selection
  \setmainfont[]{Times New Roman}
\fi
% Use upquote if available, for straight quotes in verbatim environments
\IfFileExists{upquote.sty}{\usepackage{upquote}}{}
\IfFileExists{microtype.sty}{% use microtype if available
  \usepackage[]{microtype}
  \UseMicrotypeSet[protrusion]{basicmath} % disable protrusion for tt fonts
}{}
\makeatletter
\@ifundefined{KOMAClassName}{% if non-KOMA class
  \IfFileExists{parskip.sty}{%
    \usepackage{parskip}
  }{% else
    \setlength{\parindent}{0pt}
    \setlength{\parskip}{6pt plus 2pt minus 1pt}}
}{% if KOMA class
  \KOMAoptions{parskip=half}}
\makeatother
\usepackage{color}
\usepackage{fancyvrb}
\newcommand{\VerbBar}{|}
\newcommand{\VERB}{\Verb[commandchars=\\\{\}]}
\DefineVerbatimEnvironment{Highlighting}{Verbatim}{commandchars=\\\{\}}
% Add ',fontsize=\small' for more characters per line
\newenvironment{Shaded}{}{}
\newcommand{\AlertTok}[1]{\textcolor[rgb]{1.00,0.00,0.00}{\textbf{#1}}}
\newcommand{\AnnotationTok}[1]{\textcolor[rgb]{0.38,0.63,0.69}{\textbf{\textit{#1}}}}
\newcommand{\AttributeTok}[1]{\textcolor[rgb]{0.49,0.56,0.16}{#1}}
\newcommand{\BaseNTok}[1]{\textcolor[rgb]{0.25,0.63,0.44}{#1}}
\newcommand{\BuiltInTok}[1]{\textcolor[rgb]{0.00,0.50,0.00}{#1}}
\newcommand{\CharTok}[1]{\textcolor[rgb]{0.25,0.44,0.63}{#1}}
\newcommand{\CommentTok}[1]{\textcolor[rgb]{0.38,0.63,0.69}{\textit{#1}}}
\newcommand{\CommentVarTok}[1]{\textcolor[rgb]{0.38,0.63,0.69}{\textbf{\textit{#1}}}}
\newcommand{\ConstantTok}[1]{\textcolor[rgb]{0.53,0.00,0.00}{#1}}
\newcommand{\ControlFlowTok}[1]{\textcolor[rgb]{0.00,0.44,0.13}{\textbf{#1}}}
\newcommand{\DataTypeTok}[1]{\textcolor[rgb]{0.56,0.13,0.00}{#1}}
\newcommand{\DecValTok}[1]{\textcolor[rgb]{0.25,0.63,0.44}{#1}}
\newcommand{\DocumentationTok}[1]{\textcolor[rgb]{0.73,0.13,0.13}{\textit{#1}}}
\newcommand{\ErrorTok}[1]{\textcolor[rgb]{1.00,0.00,0.00}{\textbf{#1}}}
\newcommand{\ExtensionTok}[1]{#1}
\newcommand{\FloatTok}[1]{\textcolor[rgb]{0.25,0.63,0.44}{#1}}
\newcommand{\FunctionTok}[1]{\textcolor[rgb]{0.02,0.16,0.49}{#1}}
\newcommand{\ImportTok}[1]{\textcolor[rgb]{0.00,0.50,0.00}{\textbf{#1}}}
\newcommand{\InformationTok}[1]{\textcolor[rgb]{0.38,0.63,0.69}{\textbf{\textit{#1}}}}
\newcommand{\KeywordTok}[1]{\textcolor[rgb]{0.00,0.44,0.13}{\textbf{#1}}}
\newcommand{\NormalTok}[1]{#1}
\newcommand{\OperatorTok}[1]{\textcolor[rgb]{0.40,0.40,0.40}{#1}}
\newcommand{\OtherTok}[1]{\textcolor[rgb]{0.00,0.44,0.13}{#1}}
\newcommand{\PreprocessorTok}[1]{\textcolor[rgb]{0.74,0.48,0.00}{#1}}
\newcommand{\RegionMarkerTok}[1]{#1}
\newcommand{\SpecialCharTok}[1]{\textcolor[rgb]{0.25,0.44,0.63}{#1}}
\newcommand{\SpecialStringTok}[1]{\textcolor[rgb]{0.73,0.40,0.53}{#1}}
\newcommand{\StringTok}[1]{\textcolor[rgb]{0.25,0.44,0.63}{#1}}
\newcommand{\VariableTok}[1]{\textcolor[rgb]{0.10,0.09,0.49}{#1}}
\newcommand{\VerbatimStringTok}[1]{\textcolor[rgb]{0.25,0.44,0.63}{#1}}
\newcommand{\WarningTok}[1]{\textcolor[rgb]{0.38,0.63,0.69}{\textbf{\textit{#1}}}}
\usepackage{longtable,booktabs,array}
\newcounter{none} % for unnumbered tables
\usepackage{calc} % for calculating minipage widths
% Correct order of tables after \paragraph or \subparagraph
\usepackage{etoolbox}
\makeatletter
\patchcmd\longtable{\par}{\if@noskipsec\mbox{}\fi\par}{}{}
\makeatother
% Allow footnotes in longtable head/foot
\IfFileExists{footnotehyper.sty}{\usepackage{footnotehyper}}{\usepackage{footnote}}
\makesavenoteenv{longtable}
\ifLuaTeX
\usepackage[bidi=basic,shorthands=off]{babel}
\else
\usepackage[bidi=default,shorthands=off]{babel}
\fi
\ifPDFTeX
\else
\babelfont{rm}[]{Times New Roman}
\fi
\ifLuaTeX
  \usepackage{selnolig} % disable illegal ligatures
\fi
\setlength{\emergencystretch}{3em} % prevent overfull lines
\providecommand{\tightlist}{%
  \setlength{\itemsep}{0pt}\setlength{\parskip}{0pt}}
\usepackage{fancyhdr}
\usepackage{enumitem}
\usepackage{caption}
\usepackage{bm}
\usepackage{float}
\usepackage{microtype}
\usepackage{xurl}
\usepackage{ragged2e}

\pagestyle{fancy}
\fancyhf{}
\fancyfoot[C]{\thepage}
\renewcommand{\headrulewidth}{0pt}

\setlength{\parindent}{0pt}
\setlength{\parskip}{6pt plus 2pt minus 1pt}
\setlist{nosep,leftmargin=*}
\captionsetup{font=small,labelfont=bf}
\setlength{\emergencystretch}{3em}
\sloppy
\hyphenation{im-ple-men-ta-tion op-er-a-tion-al-i-za-tion}
\usepackage{bookmark}
\IfFileExists{xurl.sty}{\usepackage{xurl}}{} % add URL line breaks if available
\urlstyle{same}
\hypersetup{
  pdftitle={Causal Fermion Systems and Super Information Theory},
  pdfauthor={Micah Blumberg},
  pdflang={en-US},
  colorlinks=true,
  linkcolor={blue},
  filecolor={Maroon},
  citecolor={blue},
  urlcolor={blue},
  pdfcreator={LaTeX via pandoc}}

\author{}
\date{}

\begin{document}

\begin{center}

{\LARGE\bfseries Causal Fermion Systems and Super Information Theory\par}

\vspace{0.35em}

{\large\itshape A Four-Lens Comparison of Effective Equations, Relational Architecture, and Public Genealogy\par}

\vspace{0.8em}

{\small\textbf{Author: Micah Blumberg}\par}

{\small Self Aware Networks Research Institute\par}

{\small\url{https://selfawarenetworks.com}\par}

{\small November 20, 2026 (Version 2 Model-Identity Appendix and Exact November Action Truncation Audit to DOI 10.5281/zenodo.21693926)\par}

\end{center}

\vspace{0.65em}

\begin{abstract}

Causal Fermion Systems (CFS) and Super Information Theory (SIT) approach gravity and information from different mathematical starting points. This paper asks two separate questions. First, do their equations meet in a controlled regime? Second, do their dated explanatory architectures implement partly equivalent operations even where the mathematical objects differ? Fischer and Finster's 2026 curved-spacetime continuum analysis provides the narrow mathematical comparison: under its regulated-continuum and constant-coupling assumptions, the CFS equations reduce to Einstein gravity coupled to one on-shell Dirac excitation. The corrected SIT Version 80 Draft 13 family contains the same equation class in a fully specified submodel with canonical scalar kinetics, a stable quadratic potential, unit matter-frame and electromagnetic functions, neutral minimally coupled Dirac matter, and a vanishing Maxwell background. CFS, the SIT submodel, and a direct Einstein-Dirac null all match three field-equation signatures and the same weak-field radial acceleration. The overlap is real but has zero discriminatory value against that null.

The historical result is more specific. Micah Blumberg's publicly frozen April 2017 Neural Lace source already joins a distributed time-space pattern, coincidence-defined information, receiver-to-receiver communication, feedback as the conscious relation, and a decoded system-level representation. Fischer and Paganini's 2025 CFS ontology later describes spacetime as a web of correlations, an emergent background, and experiments conducted from within the system. A March 2026 CFS comparison then identifies relation encoding by a generalized two-point correlator as a key CFS innovation and transfers that construction across theories. The ordered architecture passes the functional and narrative lenses at a partial, cross-domain level. Blumberg is earlier on that bounded distributed-representation / feedback architecture delta, while CFS and named physics precursors retain priority for their older operator, correlation-spacetime, causal-action, and continuum-limit machinery. Temporal nonlocality and coherence are also resolved rather than left pending: the frameworks share broad problem architecture but not a microscopic temporal kernel, stochastic collapse generator, or common coherence invariant. No operator-measure-to-field map, active-sector equivalence, finite-regulator error bound, collapse transport, historical access, or influence follows. Einstein-Dirac remains this paper's mathematical ceiling.

\end{abstract}

\section{Scope and evidence boundary}\label{scope-and-evidence-boundary}

\url{https://selfawarenetworks.com}

\section{1. Question and Claim Ceiling}\label{question-and-claim-ceiling}

The comparison question is:

\begin{quote}
Does a defined map carry specified CFS objects, dynamics, invariants, or observables into a frozen SIT model while preserving their mathematical and empirical roles?
\end{quote}

This question is narrower than whether the frameworks use similar words. It is also narrower than whether later CFS work resembles earlier ideas in Micah Blumberg's broader corpus. A mathematical relation, a narrative architecture, and a historical priority finding require different evidence.

\subsection{1.1 The problem in plain language}\label{the-problem-in-plain-language}

Two theories can produce the same familiar equation for very different reasons. Conversely, two theories can organize information through the same sequence of operations while using different objects, equations, and vocabulary. Treating the first case as a unique match would overstate the mathematics. Treating the second case as no match at all would miss a real functional and historical relation.

This paper therefore separates four tasks:

\begin{enumerate}
\def\labelenumi{\arabic{enumi}.}
\tightlist
\item
  reconstruct what each framework receives, transforms, and produces;
\item
  test whether a typed mathematical map preserves those roles;
\item
  compare the ordered explanatory architecture atom by atom; and
\item
  place each matched atom on a dated before-and-after timeline.
\end{enumerate}

Only after those operations are visible are the paper's comparison terms introduced. A \textbf{typed map} is a rule that says what mathematical object is sent to what other object and what structure is preserved. A \textbf{narrative architecture} is the ordered causal sequence that a theory says turns inputs into outputs. A \textbf{priority delta} is a bounded matched atom for which one dated public formulation is earlier; it is not a whole-theory ownership claim.

\subsection{1.2 Relation levels and claim ceiling}\label{relation-levels-and-claim-ceiling}

The audit permits five relation levels:

\begin{enumerate}
\def\labelenumi{\arabic{enumi}.}
\tightlist
\item
  shared vocabulary;
\item
  architectural analogy;
\item
  typed object or dynamics map;
\item
  proved limiting or invariant-preserving relation;
\item
  shared or competing empirical prediction.
\end{enumerate}

Level 2 is reached for one generic variational sequence. At the effective field-equation level, the frozen CFS regime, SIT submodel, and direct Einstein-Dirac null also occupy the same equation class. That result does not supply a map between primitive objects or derivations, and its zero discrimination against the null prevents promotion to a distinctive cross-framework mathematical relation. The newer-source audit additionally finds a partial Level 2 pass for the ordered distributed-representation and internal-feedback architecture. That positive narrative result is not erased by the failure of a whole-framework typed map.

\section{2. Frozen Sources}\label{frozen-sources}

\subsection{2.1 CFS introduction}\label{cfs-introduction}

Finster, Kindermann, and Treude, \emph{Causal Fermion Systems: An Introduction to Fundamental Structures, Methods and Applications}, arXiv:2411.06450v2, 2025-07-17. Local PDF SHA-256: \nolinkurl{CFC535D3BCEFEEDF59FB6DC26F42AE0673647B0FFE1C4DC4E3EF0325F976B9CA}.

The source-fixed anchors are:

\begin{itemize}
\tightlist
\item
  Section 5.4, p.~90, Definition 5.4.1: a causal fermion system is built from a separable complex Hilbert space, an operator set with a bounded signature, and a positive measure.
\item
  Equation 5.24: spacetime is \texttt{M\ :=\ supp\ rho}.
\item
  Section 5.6, pp.~94-99, equations 5.35-5.39: the spectral Lagrangian, causal action, and volume, trace, and boundedness constraints.
\item
  Definition 5.6.1: the spectral classification of operator products used for causal structure.
\item
  Chapter 7, beginning p.~133: Euler-Lagrange equations for the causal variational principle.
\item
  Chapter 9, beginning p.~157: surface-layer integrals and conservation laws.
\item
  Chapter 21, pp.~335-346, especially Sections 21.2-21.3: continuum-limit formalism and result overview.
\item
  Chapter 22, beginning p.~347: connection to quantum field theory.
\end{itemize}

\subsection{2.2 CFS effective-collapse paper}\label{cfs-effective-collapse-paper}

Finster, Kleiner, and Paganini, \emph{Causal Fermion Systems as an Effective Collapse Theory}, arXiv:2405.19254v2, \emph{Journal of Physics A} 57 (2024) 395303. Local PDF SHA-256: \nolinkurl{FF8DE2D426B316613BDC7CCA44AC25C534B25B87BDFA83356215641347239944}.

The source-fixed anchors are:

\begin{itemize}
\tightlist
\item
  Abstract, p.~1: nonlinear and stochastic corrections, a deterministic Kossakowski-Lindblad statistical-operator equation, Born-compatible dynamical collapse, probability-integral conservation, and microscopic temporal nonlocality are claims of the effective model.
\item
  Section 3.3, p.~16: Kossakowski-Lindblad dynamics.
\item
  Section 4.5, p.~21, Theorem 4.1: the stated deterministic equation for the statistical operator, with the paper's approximations and error analysis.
\end{itemize}

These results remain CFS results. They cannot be transferred to SIT by analogy.

\subsection{2.3 CFS relational-ontology paper}\label{cfs-relational-ontology-paper}

Fischer and Paganini, \emph{Causal Fermion Systems: Spacetime as the Web of Correlations of a Many-Body Quantum System}, arXiv:2504.19272v1, submitted 2025-04-27T15:05:59Z. Local PDF SHA-256: \nolinkurl{4320E599C07AED1B9D2F925E015442937426A96628A1686169D4D07455B71EC2}.

The paper describes its ontology as a first proposal for a full physical interpretation of CFS independent of the effective continuum description. It states that system information is encoded in two-point correlation strengths and position observables, describes spacetime as a web of correlations, and explains an apparent background and experiments from within the system. It also names prior correlation and relational lineages, including Rovelli, Kempf, and Reitz, Soda, and Kempf. Those precursors remain part of the genealogy.

\subsection{2.4 CFS generalized-correlator comparison}\label{cfs-generalized-correlator-comparison}

Farnsworth, Finster, Paganini, and Singh, \emph{Causal Fermion Systems, Non-Commutative Geometry and Generalized Trace Dynamics}, arXiv:2603.05018v2. Version 1 was submitted 2026-03-05T10:05:31Z and version 2 on 2026-03-10T11:06:15Z. Local PDF SHA-256: \nolinkurl{2B2D212F585673755F2EC357C42A4F88D6CDE0452C1544E6B0C70AB6C230C711}.

The abstract and Sections 3.1-3.4 identify the encoding of relations between spacetime points by a generalized two-point correlator as a key CFS innovation, compare the relevant objects and action principles across three frameworks, and show how the correlator construction can be transferred. The source strengthens the relational comparison but does not create a map to SIT.

\subsection{2.5 CFS curved-spacetime continuum result}\label{cfs-curved-spacetime-continuum-result}

Fischer and Finster, \emph{The Continuum Limit Analysis of Causal Fermion Systems for Curved Spacetimes}, arXiv:2605.30199v1, submitted 2026-05-28. The locally frozen PDF has SHA-256 \nolinkurl{83C15821775F82B83280100BE8B1F09206B3173AEFB0FD7417691DE87DA7DC54}.

The source direction is essential. It begins with a globally hyperbolic spacetime, constructs a corresponding regulated causal fermion system, and analyzes restricted Euler-Lagrange equations as the regulator tends to zero. Theorem 5.2 and Corollary 5.3 recover the vacuum Einstein equation. Theorem 5.4 and Corollary 5.5 recover the Einstein equation coupled to an on-shell Dirac excitation when the effective gravitational coupling is constant. This is not a theorem that an arbitrary causal fermion system reconstructs an arbitrary smooth spacetime.

\subsection{2.6 SIT Version 80}\label{sit-version-80}

The SIT comparison target is the governed Version 80 Draft 13 current manuscript at \nolinkurl{super-information-theory-v80.md}, SHA-256 \nolinkurl{9DA6276C4B3243E5367C14266C9191CBCB576F806362C3D4CA5B56A4673550F8}.

Sections 4-8 define the relevant source ceiling: an Einstein-frame metric, neutral scalar fields \texttt{tau} and \texttt{sigma}, an independent Maxwell connection, matter fields, a conformal physical metric, and a local two-derivative action. Draft 13 also makes coherence measures protocol- and basis-specific, requires a calibrated state-to-field kernel before a measured preparation statistic is treated as a field, and preserves the public 2022 QGTCD genealogy. It does not supply a controlled quantum EFT, a stochastic collapse generator, a microscopic temporal-nonlocality kernel, or a derivation of its spacetime manifold from a more primitive operator measure.

\section{3. Independent Reconstructions}\label{independent-reconstructions}

\subsection{3.1 CFS object-action-output chain}\label{cfs-object-action-output-chain}

At the resolution used here, CFS has the typed chain

\begin{Shaded}
\begin{Highlighting}[]
\NormalTok{(H, F, rho; constraints)}
\NormalTok{    {-}\textgreater{} spectral data of xy}
\NormalTok{    {-}\textgreater{} Lagrangian L(x,y)}
\NormalTok{    {-}\textgreater{} double{-}integral causal action S(rho)}
\NormalTok{    {-}\textgreater{} variation of rho under constraints}
\NormalTok{    {-}\textgreater{} Euler{-}Lagrange conditions}
\NormalTok{    {-}\textgreater{} continuum{-}limit effective structures under additional assumptions}
\end{Highlighting}
\end{Shaded}

The primitive configuration is a measure on an operator space. Spacetime is identified with the support of that measure. Causal relations are not inserted as the causal cones of an assumed Lorentzian metric; they arise from spectral properties of products of point operators. Familiar local field descriptions are sought through continuum-limit constructions.

\subsection{3.2 SIT v80 object-action-output chain}\label{sit-v80-object-action-output-chain}

The corrected SIT target has the chain

\begin{Shaded}
\begin{Highlighting}[]
\NormalTok{(M, g, tau, sigma, A, chi; functions and boundary data)}
\NormalTok{    {-}\textgreater{} local two{-}derivative spacetime action}
\NormalTok{    {-}\textgreater{} variation of g, tau, sigma, and A}
\NormalTok{    {-}\textgreater{} field equations and conservation identities}
\NormalTok{    {-}\textgreater{} physical{-}frame observables if a viable benchmark is fixed}
\end{Highlighting}
\end{Shaded}

The differentiable manifold and metric sector are inputs. The matter metric is conformally related to the Einstein-frame metric. The scalar, Maxwell, and matter sectors have familiar field-theoretic types. Version 80 is corrective: it constrains what may be claimed but does not derive a complete fundamental physics theory.

\section{4. The Surviving Functional Analogy}\label{the-surviving-functional-analogy}

Both systems instantiate the directed pattern

\begin{Shaded}
\begin{Highlighting}[]
\NormalTok{admissible configuration {-}\textgreater{} variational functional {-}\textgreater{} stationarity equations}
\end{Highlighting}
\end{Shaded}

This is a candidate functional analogy. It is useful because it identifies where a translation might be attempted. It is weak because variational architecture is common across mathematical physics. The shared sequence does not imply that the configurations, actions, constraints, equations, or observables are equivalent.

The implementations differ at the surface immediately:

\begin{itemize}
\tightlist
\item
  CFS varies a positive measure on a space of finite-rank operators; SIT v80 varies local fields on an assumed manifold.
\item
  The CFS causal action is a nonlocal double integral over operator pairs; the displayed SIT action is local and two-derivative.
\item
  CFS geometry and causal relations are reconstructed from operator-measure data; SIT v80 begins with metric geometry.
\item
  CFS effective-collapse results use a separate approximation stack; SIT v80 contains no collapse dynamics.
\end{itemize}

\section{5. Typed Candidate Map}\label{typed-candidate-map}

Let

\begin{Shaded}
\begin{Highlighting}[]
\NormalTok{C\_CFS = \{rho in positive measures on F satisfying C\_volume, C\_trace, C\_bound\}}

\NormalTok{C\_SIT(M) = \{(g,tau,sigma,A,chi; V,Omega,f\_EM,Z\_tau,Z\_sigma) admissible on M\}}
\end{Highlighting}
\end{Shaded}

A forward comparison map would have the type

\begin{Shaded}
\begin{Highlighting}[]
\NormalTok{T : C\_CFS {-}\textgreater{} disjoint\_union over M of C\_SIT(M).}
\end{Highlighting}
\end{Shaded}

It must therefore do more than send the symbol \texttt{rho} to a scalar called \texttt{rho}. Given a CFS measure, \texttt{T} must:

\begin{enumerate}
\def\labelenumi{\arabic{enumi}.}
\tightlist
\item
  select or derive a differentiable manifold from \texttt{supp\ rho};
\item
  derive a Lorentzian metric and physical matter metric;
\item
  derive the two scalar fields and their kinetic functions;
\item
  derive the Maxwell connection, gauge normalization, and charges;
\item
  derive the matter content and its action;
\item
  map the CFS constraints and causal action to the SIT action or its limit;
\item
  preserve at least one specified invariant or observable.
\end{enumerate}

No rule satisfying steps 1-7 is present in the frozen sources.

\subsection{5.1 Domain failure}\label{domain-failure}

The support of a measure on operator space need not have a smooth manifold structure. SIT v80 presupposes one. Consequently \texttt{T} is not defined on the full CFS configuration space unless its domain is restricted to a controlled class with a proved continuum construction.

\subsection{5.2 Field-recovery failure}\label{field-recovery-failure}

Even after a continuum spacetime is available, the frozen comparison supplies no unique extraction of \texttt{tau}, \texttt{sigma}, \texttt{A}, \texttt{chi}, or the model functions. The map is underdetermined.

\subsection{5.3 Dynamics failure}\label{dynamics-failure}

A common variational arrow does not establish

\begin{Shaded}
\begin{Highlighting}[]
\NormalTok{T(EL\_CFS(rho)) = EL\_SIT(T(rho)).}
\end{Highlighting}
\end{Shaded}

No intertwining relation, controlled error term, or matching boundary class is proved.

\subsection{5.4 Invariant failure}\label{invariant-failure}

CFS surface-layer conservation laws, collapse-model probability integrals, and SIT stress-energy or gauge-current identities have different definitions and domains. No shared invariant \texttt{I} has been exhibited with \texttt{I\_CFS(rho)\ =\ I\_SIT(T(rho))}.

\subsection{5.5 Reverse-reconstruction failure}\label{reverse-reconstruction-failure}

A reverse map from the finite SIT field tuple to a universal measure on operator space is also absent. Information loss, nonuniqueness, regularization, and gauge equivalence have not been controlled.

The typed-map result is therefore \nolinkurl{failed-at-current-source-ceiling}. This is not a proof that no map exists. It identifies the missing construction.

\section{6. Collapse, Nonlocality, and Coherence}\label{collapse-nonlocality-and-coherence}

The effective-collapse paper establishes claims inside its own model and approximation regime. Its stochastic corrections, Kossakowski-Lindblad equation, probability-integral conservation, and microscopic temporal nonlocality are a joined CFS mechanism. SIT Draft 13 does not contain the same stochastic generator, quantum state space, collapse observable, or microscopic temporal kernel. Its baseline is a local two-derivative scalar-tensor model. The public 2022 QGTCD genealogy instead asks whether locally varying temporal opportunities bias motion, while the separate SuperTimePosition program asks about temporal placement and phase relation. These are temporally organized problems, but they are not the CFS collapse kernel.

The temporal-nonlocality comparison is therefore closed as:

\texttt{partial\ narrative\ relation;\ technical\ non-equivalence;\ CFS-specific\ collapse\ result\ retained}.

The comparison must also keep four meanings apart:

\begin{enumerate}
\def\labelenumi{\arabic{enumi}.}
\tightlist
\item
  statistical dependence;
\item
  neural phase coherence;
\item
  computational integration;
\item
  quantum coherence.
\end{enumerate}

Only the fourth is directly relevant to a quantum-collapse comparison, and even there the exact resource, density operator, basis, dynamics, and measurement protocol must be fixed. Draft 13 makes that boundary operational: a coherence statistic must declare its basis and preparation, and any state-to-field use requires a calibrated kernel with units, support, uncertainty, and transfer tests. CFS probability-integral conservation, CFS surface-layer conservation, and SIT stress-energy or gauge-current identities remain different invariants.

The coherence comparison is therefore also closed:

\texttt{no\ shared\ coherence-conservation\ theorem;\ partial\ architecture\ only;\ no\ unresolved\ invariant\ claim}.

\section{7. Dated Atom Genealogy and Before-and-After Comparison}\label{dated-atom-genealogy-and-before-and-after-comparison}

\subsection{7.1 Micah's public 2017 architecture}\label{micahs-public-2017-architecture}

Three Neural Lace Podcast recordings were published by Silicon Valley Global News on 11, 13, and 22 April 2017. Their public metadata, transcript bytes, SHA-256 hashes, and available Wayback CDX captures are frozen in the package receipt {[}8-10{]}. The first episode supplies the load-bearing sequence:

\begin{enumerate}
\def\labelenumi{\arabic{enumi}.}
\tightlist
\item
  lines 165-189 describe a changing three-dimensional time-space ionic pattern as a distributed representation of a sensory object;
\item
  lines 335-371 ask whether active conscious observation is distinguished by a returned feedback relation rather than broadcast-only transmission;
\item
  lines 421-457 propose an alternating receive / transmit protocol specialized by information type; and
\item
  lines 471-495 join distributed concept components, coincidence detection, bit-like information, and a transmitted pattern.
\end{enumerate}

This sequence is a public ancestor of Micah's later receiver-relative, distributed-representation, and internal-observer architecture. It is not the later SIT field equations, a CFS universal measure, or a generalized two-point-correlator theorem.

\subsection{7.2 Micah's public 2022-2026 physics sequence}\label{micahs-public-2022-2026-physics-sequence}

The current SIT atlas closes the paper-relevant later genealogy:

{\def\LTcaptype{none} % do not increment counter
\begingroup\footnotesize\setlength{\tabcolsep}{4pt}
\begin{longtable}[]{@{}
  >{\RaggedRight\arraybackslash}p{(\linewidth - 4\tabcolsep) * \real{0.3333}}
  >{\RaggedRight\arraybackslash}p{(\linewidth - 4\tabcolsep) * \real{0.3333}}
  >{\RaggedRight\arraybackslash}p{(\linewidth - 4\tabcolsep) * \real{0.3333}}@{}}
\toprule\noalign{}
\begin{minipage}[b]{\linewidth}\RaggedRight
Public date
\end{minipage} & \begin{minipage}[b]{\linewidth}\RaggedRight
Source
\end{minipage} & \begin{minipage}[b]{\linewidth}\RaggedRight
Bounded atom
\end{minipage} \\
\midrule\noalign{}
\endhead
\bottomrule\noalign{}
\endlastfoot
2022-06-10 & Git commit \nolinkurl{40c0b37fc9376a6630ad9d3b09d911997c1c17f4}, \texttt{a0010z.md} {[}11{]} & oscillators, mass, density, scale, phase structure, curvature, and time dilation joined in one exploratory record \\
2022-06-12 & Git commit \nolinkurl{4c451f016cd076bcb635f4305a55efe54bbf4870}, \texttt{a0387z.md} {[}12{]} & a gradient of temporal opportunities around mass is proposed to bias directional motion \\
2022-11-28 & Git commit \nolinkurl{4e0cf1ee0975c01dc57182368497d3f0cd864030}, \texttt{a0322z.md} {[}13{]} & QGTCD questions extend to rotation curves, MOND, Tully-Fisher, dark-sector, and Hubble-tension consequences \\
2024-2025 & named QGTCD, Dark Time, Super Dark Time, and SIT publications & the component questions become named field and synthesis programs \\
2026-07-26 & SIT Version 80 Draft 13 {[}6{]} & a corrected scalar-tensor scaffold, state-to-field calibration requirement, conditional gradient-to-motion operator, and explicit public genealogy \\
\end{longtable}
\endgroup
}

The later stages increase mathematical and experimental specificity; they do not replace the 2022 public atoms or assign Draft 13 equations to 2022.

\subsection{7.3 CFS before and after}\label{cfs-before-and-after}

The CFS chronology must also be versioned:

{\def\LTcaptype{none} % do not increment counter
\begingroup\footnotesize\setlength{\tabcolsep}{4pt}
\begin{longtable}[]{@{}
  >{\RaggedRight\arraybackslash}p{(\linewidth - 2\tabcolsep) * \real{0.5000}}
  >{\RaggedRight\arraybackslash}p{(\linewidth - 2\tabcolsep) * \real{0.5000}}@{}}
\toprule\noalign{}
\begin{minipage}[b]{\linewidth}\RaggedRight
Stage
\end{minipage} & \begin{minipage}[b]{\linewidth}\RaggedRight
Retained result
\end{minipage} \\
\midrule\noalign{}
\endhead
\bottomrule\noalign{}
\endlastfoot
foundational CFS and the 2024/2025 introduction {[}1{]} & operator space, universal measure, spectral causal action, support as spacetime, operator-product causal relations, and continuum-limit program are earlier CFS machinery \\
2024 effective-collapse paper {[}2{]} & an exact, approximation-bounded CFS collapse mechanism with microscopic temporal nonlocality \\
April 2025 correlation-web paper {[}3{]} & a proposed full physical interpretation adds information as two-point correlations and position observables, spacetime as a web of correlations, an emergent background, and experiments from within \\
March 2026 cross-theory comparison {[}4{]} & relation encoding by a generalized two-point correlator is foregrounded and transferred across CFS, NCG, and GTD \\
May 2026 curved-spacetime result {[}5{]} & a restricted construction recovers Einstein-Dirac equations and supplies this paper's mathematical benchmark \\
\end{longtable}
\endgroup
}

The 2025 source explicitly credits earlier relational and correlation-spacetime lineages. In particular, it names Rovelli, Kempf's 2021 replacement of distance by correlation, and Reitz, Soda, and Kempf's 2023 spacetime-emergence model. Those are genuine precursors to the physics-specific ontology. The Micah result concerns the earlier cross-domain ordered architecture, not ownership of all correlation-based physics.

\subsection{7.4 Atom-by-atom result}\label{atom-by-atom-result}

{\def\LTcaptype{none} % do not increment counter
\begingroup\footnotesize\setlength{\tabcolsep}{4pt}
\begin{longtable}[]{@{}
  >{\RaggedRight\arraybackslash}p{(\linewidth - 4\tabcolsep) * \real{0.3333}}
  >{\RaggedRight\arraybackslash}p{(\linewidth - 4\tabcolsep) * \real{0.3333}}
  >{\RaggedRight\arraybackslash}p{(\linewidth - 4\tabcolsep) * \real{0.3333}}@{}}
\toprule\noalign{}
\begin{minipage}[b]{\linewidth}\RaggedRight
Atom
\end{minipage} & \begin{minipage}[b]{\linewidth}\RaggedRight
Match
\end{minipage} & \begin{minipage}[b]{\linewidth}\RaggedRight
Dated disposition
\end{minipage} \\
\midrule\noalign{}
\endhead
\bottomrule\noalign{}
\endlastfoot
distributed local states become a system-level representation & partial functional pass & Micah's public neural articulation is earlier than the 2025 CFS ontology articulation \\
local relations or coincidences carry system-level information & partial functional and narrative pass & Micah 2017 is earlier on the bounded neural implementation; older CFS and named physics precedents retain their own mathematical priority \\
an apparent background or representation is constructed from distributed relations & partial narrative pass & later convergence in formulation; domain and object types differ \\
the system samples or acts from within the constructed pattern and receives returned information & narrative-architecture pass & Micah's 2017 feedback architecture is earlier than the 2025 experiment-from-within exposition \\
Einstein-Dirac field equations and weak-field acceleration & exact effective-class pass & real overlap, but zero discrimination against the direct null \\
temporal nonlocality and collapse & broad problem relation only & CFS's 2024 technical mechanism is retained; no SIT kernel equivalence \\
coherence conservation & no shared invariant & meanings are operationally non-commensurate \\
\end{longtable}
\endgroup
}

\section{8. Four-Lens Audit}\label{four-lens-audit}

\subsection{8.1 Equivalence lens}\label{equivalence-lens}

Two bounded passes survive:

\begin{enumerate}
\def\labelenumi{\arabic{enumi}.}
\tightlist
\item
  a generic action-to-stationarity schema, which the negative controls show is nondiscriminating; and
\item
  a partial functional map from distributed local relations through aggregation to a system-level representation used from within the system.
\end{enumerate}

The Einstein-Dirac equation-class match is exact in the frozen regime. Primitive-object, microscopic-dynamics, invariant, collapse, and whole-framework equivalence do not follow. No bidirectional functor or isomorphism is established.

\subsection{8.2 Narrative and architecture lens}\label{narrative-and-architecture-lens}

The atomized shared architecture is:

\begin{Shaded}
\begin{Highlighting}[]
\NormalTok{distributed local states}
\NormalTok{    {-}\textgreater{} relations or coincidences among them}
\NormalTok{    {-}\textgreater{} aggregated system{-}level pattern}
\NormalTok{    {-}\textgreater{} effective background or representation}
\NormalTok{    {-}\textgreater{} system samples or acts from within that pattern}
\NormalTok{    {-}\textgreater{} returned information changes the next state}
\end{Highlighting}
\end{Shaded}

This ordered sequence passes at a partial cross-domain level. Differences in mathematical carrier and physical domain do not erase the functional order. They bound the result to narrative and operational architecture rather than a full mathematical identity.

\subsection{8.3 Detective and chronology lens}\label{detective-and-chronology-lens}

Starting from the later CFS conclusion - an internally observed relational system whose apparent background is reconstructed from distributed correlations - the audit works backward to the public 2017 Neural Lace sequence. That source contains the earlier ordered architecture. It also works backward through CFS and the named correlation-physics lineages, preventing a false claim that their technical machinery began in 2025. The resulting disposition is \nolinkurl{Micah-earlier-on-the-bounded-cross-domain-architecture-delta}, not a claim of access, copying, or misconduct.

\subsection{8.4 Longitudinal signature-cluster lens}\label{longitudinal-signature-cluster-lens}

The older CFS record already has relational mathematics, but its later presentation changes in a recognizable sequence. The April 2025 source calls its correlation-web account a first full physical interpretation, emphasizes information and the emergent observer-relative background, and describes experiments from within. The March 2026 source then identifies generalized relation encoding as a key innovation and exports the construction across frameworks. That later cluster is closer in narrative form to Micah's 2017 distributed representation, coincidence, feedback, and internally sampled process than the earlier CFS presentation was.

This lens establishes later convergence in formulation. The available record does not establish why the convergence occurred.

\subsection{8.5 Combined disposition}\label{combined-disposition}

The four lenses do not yield one global yes/no verdict. The correct atomized result is:

\begin{itemize}
\tightlist
\item
  \nolinkurl{candidate-functional-analogy} for the generic variational sequence;
\item
  \nolinkurl{failed-at-current-source-ceiling} for the tested typed map;
\item
  \nolinkurl{partial-functional-and-narrative-match} for the ordered distributed relation-to-representation architecture;
\item
  \nolinkurl{Micah-earlier-on-bounded-cross-domain-architecture-delta} for the public 2017-to-2025 chronology;
\item
  \nolinkurl{CFS-and-named-priors-retained} for their older mathematical and physics-specific machinery;
\item
  \nolinkurl{technical-non-equivalence} for temporal nonlocality and collapse;
\item
  \nolinkurl{no-shared-invariant} for coherence conservation; and
\item
  \nolinkurl{no-influence-evidence} for historical access.
\end{itemize}

\section{9. Controlled Calibration Against Variational Baselines}\label{controlled-calibration-against-variational-baselines}

\subsection{9.1 Why a control set is required}\label{why-a-control-set-is-required}

The sequence

\begin{Shaded}
\begin{Highlighting}[]
\NormalTok{admissible configuration {-}\textgreater{} variational functional {-}\textgreater{} stationarity equations}
\end{Highlighting}
\end{Shaded}

is not distinctive merely because it is present in both target frameworks. Classical mechanics, Maxwell theory, Yang-Mills theory, general relativity, scalar-tensor-Maxwell models, and free-energy field models also instantiate that sequence. A structural comparison needs a negative control: a feature that matches every ordinary variational framework cannot carry much weight as evidence of a special CFS / SIT relationship.

\subsection{9.2 Frozen slots and controls}\label{frozen-slots-and-controls}

The calibration model codes eight source-facing slots:

\begin{enumerate}
\def\labelenumi{\arabic{enumi}.}
\tightlist
\item
  primitive carrier;
\item
  base domain;
\item
  action architecture;
\item
  variation target;
\item
  status of geometry;
\item
  status of gauge structure;
\item
  invariant family; and
\item
  empirical-output status.
\end{enumerate}

The frozen controls are general relativity, scalar-tensor-Maxwell theory, Maxwell electrodynamics, Yang-Mills theory, classical mechanics, and a Cahn-Hilliard-type free-energy model. The categories deliberately encode only the comparison granularity used in this paper. They do not represent complete theories. Exact equality in a slot is a conservative test; semantic similarity without a typed map is not counted.

The executable specification is \nolinkurl{model/run_draft3_typed_signature_controls.py}; its deterministic output is \nolinkurl{results/draft3-typed-signature-control-results.json}.

\subsection{9.3 Result}\label{result}

All 28 framework pairs match the coarse variational chain. Within this control set, the chain therefore has a discrimination fraction of zero. The CFS / SIT pair has zero exact matches across the eight typed slots. CFS has zero exact typed matches with five controls and one incidental match with Yang-Mills at the broad \texttt{model\_specific} empirical-output label. SIT v80 has six exact matches with scalar-tensor-Maxwell theory, four with general relativity, three with Maxwell theory, two with Yang-Mills, and zero with the remaining two controls.

The result changes the language of the surviving analogy. It is better called a \textbf{generic variational schema} than a candidate functional isomorphism. The source-defined frameworks may still admit a partial correspondence in a controlled regime, but the three-stage chain alone supplies no discriminatory evidence for it.

\subsection{9.4 Robustness and limits}\label{robustness-and-limits}

The model is deterministic and its output is byte-reproducible. Its strongest limit is also explicit: the slot vocabulary is an analyst-coded ontology. A different but defensible ontology can change individual counts. That is why the paper publishes every category and mismatch rather than reporting one opaque similarity score. The calibration is not a continuum-limit calculation, a physical simulation, a theorem about all CFS or SIT formulations, or evidence about historical access.

\section{10. Controlled Einstein-Dirac Overlap}\label{controlled-einstein-dirac-overlap}

\subsection{10.1 Why the benchmark changed}\label{why-the-benchmark-changed}

Draft 3 proposed a joint metric / gauge calculation but did not have a concrete shared output. The 2026 Fischer-Finster result supplies a sharper target: Einstein gravity coupled to one on-shell Dirac excitation. It does not supply the SIT scalar or Maxwell sectors, so Draft 4 restricts the benchmark to the equations actually present on both sides.

\subsection{10.2 CFS regime}\label{cfs-regime}

The CFS side is the construction and limit stated in arXiv:2605.30199v1:

\begin{enumerate}
\def\labelenumi{\arabic{enumi}.}
\tightlist
\item
  a globally hyperbolic spacetime;
\item
  a quasi-free Hadamard state and corresponding fermionic projector;
\item
  chart-independent \texttt{i\ epsilon} regularization;
\item
  the continuum limit \texttt{epsilon\ -\textgreater{}\ 0};
\item
  one on-shell Dirac excitation; and
\item
  a locally translation-invariant regularization giving constant \texttt{kappa}.
\end{enumerate}

The result is

\begin{Shaded}
\begin{Highlighting}[]
\NormalTok{G\_mn + Lambda g\_mn = kappa T\_mn\^{}(Dirac),}
\NormalTok{Dirac[u,g] = 0.}
\end{Highlighting}
\end{Shaded}

The cosmological constant follows as an integration constant. The effective gravitational coupling depends on the regularization scheme.

\subsection{10.3 Frozen SIT submodel}\label{frozen-sit-submodel}

Inside the Version 80 action, choose

\begin{Shaded}
\begin{Highlighting}[]
\NormalTok{Z\_tau = Z\_sigma = 1,}
\NormalTok{Omega(tau,sigma) = 1,}
\NormalTok{f\_EM(tau,sigma) = 1,}
\NormalTok{V = M\_Pl\^{}2 Lambda}
\NormalTok{    + m\_tau\^{}2 (tau{-}tau0)\^{}2/2}
\NormalTok{    + m\_sigma\^{}2 (sigma{-}sigma0)\^{}2/2,}
\NormalTok{m\_tau\^{}2 \textgreater{} 0,  m\_sigma\^{}2 \textgreater{} 0,}
\NormalTok{tau = tau0,  sigma = sigma0,}
\NormalTok{A\_m = 0,  F\_mn = 0,}
\NormalTok{S\_m = the neutral minimally coupled Dirac action of mass m\_chi,}
\NormalTok{kappa = 1/M\_Pl\^{}2.}
\end{Highlighting}
\end{Shaded}

The scalar kinetic terms vanish at the constant background, the scalar equations are satisfied at the stationary point, the Maxwell equation is satisfied by a neutral source and vanishing field, and the matter and Einstein metrics coincide. Variation leaves the same Einstein-Dirac equation class. This is a nonempty benchmark point. It is also a standard general-relativistic submodel that turns off every proposed distinctive SIT sector.

\subsection{10.4 Partial map}\label{partial-map}

The controlled output map is

\begin{Shaded}
\begin{Highlighting}[]
\NormalTok{g\_CFS       {-}\textgreater{} g\_SIT,}
\NormalTok{u\_CFS       {-}\textgreater{} chi\_SIT,}
\NormalTok{kappa\_CFS   {-}\textgreater{} 1/M\_Pl\^{}2,}
\NormalTok{Lambda\_CFS  {-}\textgreater{} V(tau0,sigma0)/M\_Pl\^{}2.}
\end{Highlighting}
\end{Shaded}

This map is defined only after both sides have reached an effective smooth spacetime description. It does not map the CFS universal measure, operator support, regularization data, or surface-layer integrals to the complete SIT configuration.

\subsection{10.5 Common weak-field observable}\label{common-weak-field-observable}

Set \texttt{Lambda=0} and impose asymptotic flatness, a static weak field, and a localized nonrelativistic source of mass \texttt{M}. With \texttt{g\_00=-(1+2\ Phi)}, both equation sets imply

\begin{Shaded}
\begin{Highlighting}[]
\NormalTok{laplacian(Phi) = 4 pi G rho,}
\NormalTok{Phi(r) = {-}G M/r              outside the source,}
\NormalTok{a\_r = {-}G M/r\^{}2.}
\end{Highlighting}
\end{Shaded}

The radial acceleration is the common observable. Agreement is exact at this stated leading order after coupling identification. It has zero discriminatory value against general relativity because the observable follows from the shared Einstein equation.

\subsection{10.6 Executable gate audit}\label{executable-gate-audit}

The deterministic specification is \nolinkurl{model/run_draft4_einstein_dirac_overlap.py}; its output is \nolinkurl{results/draft4-einstein-dirac-overlap-results.json}. The audit reports three shared field-equation signatures out of three, one shared derived weak-field observable, and one accepted variant out of six. The five rejected variants activate a spacetime-dependent CFS coupling, an SIT scalar, an SIT Maxwell field, a nontrivial matter-frame factor, or mismatched couplings.

This audit checks source-coded eligibility. It is not a numerical physical simulation, a symbolic derivation of either source theorem, or an independent proof of the CFS continuum limit.

\subsection{10.7 Direct Einstein-Dirac null calibration}\label{direct-einstein-dirac-null-calibration}

An overlap is informative only if it distinguishes the target pair from an ordinary theory that already has the same equations. Draft 5 therefore adds the direct Einstein-Dirac theory as a null control.

The executable specification \nolinkurl{model/run_draft5_einstein_dirac_null_control.py} freezes the three common field-equation signatures and four origin slots:

\begin{enumerate}
\def\labelenumi{\arabic{enumi}.}
\tightlist
\item
  metric origin;
\item
  gravitational-coupling origin;
\item
  cosmological-constant origin; and
\item
  matter origin.
\end{enumerate}

The deterministic result is:

{\def\LTcaptype{none} % do not increment counter
\begingroup\footnotesize\setlength{\tabcolsep}{4pt}
\begin{longtable}[]{@{}lrr@{}}
\toprule\noalign{}
Pair & Shared field equations & Exact origin-slot matches \\
\midrule\noalign{}
\endhead
\bottomrule\noalign{}
\endlastfoot
CFS effective regime / SIT v80 submodel & 3/3 & 0/4 \\
CFS effective regime / Einstein-Dirac null & 3/3 & 0/4 \\
SIT v80 submodel / Einstein-Dirac null & 3/3 & 0/4 \\
\end{longtable}
\endgroup
}

The weak-field radial acceleration is also shared by all three. The discrimination fraction of the CFS / SIT equation overlap against this null is therefore zero.

This calibration does not erase the positive result. It changes its interpretation. The CFS and SIT benchmark descriptions occupy the same effective Einstein-Dirac class, but the common class is not evidence of a special relationship between them. Their frozen derivational origins remain different in every tested slot.

The origin-slot comparison is categorical and source-coded. A \texttt{0/4} score means that none of the four published labels is exactly identical at this resolution. It is not a theorem of mathematical inequivalence, and another defensible ontology could group some origins differently. The paper therefore publishes the slots and labels rather than reducing them to an opaque score.

\subsection{10.8 Regime and error ledger}\label{regime-and-error-ledger}

The equality is exact only at the stated effective equation level:

\begin{itemize}
\tightlist
\item
  CFS: restricted linearized field equations in the \texttt{epsilon\ -\textgreater{}\ 0} continuum regime with constant \texttt{kappa};
\item
  SIT: the frozen classical Einstein-Dirac submodel; and
\item
  null: classical Einstein-Dirac theory.
\end{itemize}

No finite-\texttt{epsilon} CFS error bound is transported into the comparison. No error bound covers active SIT scalar, conformal, or electromagnetic sectors. No strong-field model comparison is performed. The Newtonian observable is a standard consequence inferred from the shared Einstein equation rather than a numerical result reported by Fischer and Finster.

\subsection{10.9 Residual failures}\label{residual-failures}

No tested map covers:

\begin{itemize}
\tightlist
\item
  the universal measure \texttt{rho} to the full SIT field tuple;
\item
  CFS operator support to a derivation of the SIT manifold;
\item
  CFS regularization microstructure to \texttt{tau} or \texttt{sigma};
\item
  a CFS gauge sector to the SIT Maxwell connection;
\item
  CFS surface-layer integrals to SIT invariants;
\item
  CFS collapse dynamics to SIT dynamics; or
\item
  the spacetime-dependent \texttt{kappa(x)} case.
\end{itemize}

The constants also have different origins. In the CFS source, \texttt{Lambda} is an integration constant and \texttt{kappa} is regularization-dependent. In the frozen SIT submodel, their values are selected through \texttt{V} and \texttt{M\_Pl}. Equality of the resulting field equations does not identify those derivations.

\section{11. Versioned Micah Genealogy}\label{versioned-micah-genealogy}

The Draft 2 chronology compressed several evidentiary acts into one date. Draft 6 separated source-internal authorship or recording dates, first verified Git fixation, and first verified public availability. Draft 7 retains that ledger at \nolinkurl{MICAH-GENEALOGY-LEDGER-DRAFT6.md} and adds the current-source audit at \nolinkurl{CFS-NEWEST-SOURCE-AND-FOUR-LENS-AUDIT-DRAFT7.md}.

Seven files are now locally hash-frozen. \texttt{a0010z.md} and \texttt{b0020y.md} have first verified Git fixations on 8 June 2022. \texttt{QGTCDnote.md} is fixed in February 2024, \texttt{rawnote06.md} in October 2024, and the larger QGTCD compilation files in January 2025. These dates establish repository provenance for the corresponding bytes or ancestors. They do not show when every underlying thought originated, when each file first became public, or whether any CFS author encountered it.

This distinction matters especially for the 2025 compilations: their late Git dates do not erase earlier donor notes, but neither can those compilations be backdated without an atom-level source chain. The bounded 2017 Neural Lace row is now public-record frozen: the source channel, exact upload timestamps, recording identities, transcript hashes, passage locators, and nonzero archive capture counts are preserved in the package. This closes only that ancestry gate. It does not backdate QGTCD or SIT mathematics and does not establish access, influence, plagiarism, or derivation by CFS authors.

The governed SIT Draft 13 atlas supplies stronger public locators for the three physics atoms retained in this paper: the 10 June 2022 oscillator-mass-time record, the 12 June 2022 temporal-gradient-to-motion record, and the 28 November 2022 astrophysical-extension record. Each is cited by immutable commit and exact file. This closes the later genealogy needed for the current comparison without converting private dates or later compilations into earlier public evidence.

\section{12. Frozen Mathematical Ceiling and Separate Extensions}\label{frozen-mathematical-ceiling-and-separate-extensions}

The common-observable and null-calibration gates are closed at the effective Einstein-Dirac level. This paper stops there for a source-controlled reason. The selected CFS curved-spacetime result derives an Einstein-Dirac regime but does not furnish a mutually instantiated Maxwell benchmark. The frozen SIT submodel sets its Maxwell field to zero and its scalar to a constant vacuum. Adding a CFS gauge construction or activating SIT gauge / scalar dynamics now would select new models, parameters, and observables after the null result was known. It would therefore be a new comparison rather than a robustness check of the frozen one.

The following are preregistered extensions, not release gates for this bounded paper:

\begin{enumerate}
\def\labelenumi{\arabic{enumi}.}
\tightlist
\item
  derive or rule out a map from a specified CFS measure and regularization to the active SIT scalar and conformal sectors;
\item
  compare one explicit CFS gauge construction with the SIT Maxwell sector, including gauge group, normalization, charges, currents, and observables;
\item
  test the spacetime-dependent \texttt{kappa(x)} regime against an explicit scalar-tensor SIT benchmark;
\item
  identify and transport one invariant rather than matching equations only;
\item
  preregister a direct test of the active SIT state-to-field kernel if that kernel is specified in a future version.
\end{enumerate}

The present benchmark cannot transport CFS collapse results or establish a whole-framework equivalence. Independent CFS and mathematical-physics review would be valuable after this locally complete preprint, but it is not remaining local manuscript work and does not alter the frozen mathematical scope.

\section{13. Conclusion}\label{conclusion}

The comparison separates four results that should not be collapsed into one. First, the configuration-action-stationarity sequence is a generic variational schema and has zero discrimination in the control set. Second, the 2026 CFS continuum construction and a frozen SIT v80 Draft 13 submodel share the Einstein-Dirac effective equation class and its standard weak-field limit. Third, the direct Einstein-Dirac null shares exactly the same equations and observable. Fourth, the newer CFS sources disclose a separate, partial functional and narrative match: distributed relations are aggregated into an effective representation or background that the system uses from within.

The positive overlap is therefore real but nonspecific. It locates both target descriptions in a familiar effective class; it does not distinguish CFS from SIT, validate a distinctive SIT mechanism, or translate their microscopic objects. The metric, coupling, cosmological constant, and matter terms have different frozen origins in every pairwise comparison. Seven structures remain unmapped, including the CFS universal measure, active SIT scalar and gauge sectors, surface-layer invariants, and collapse dynamics.

The public chronology adds a bounded priority result. Micah's April 2017 Neural Lace source already contains the ordered distributed-representation, coincidence, feedback, and internal-observer architecture. The 2025 CFS ontology and March 2026 generalized-correlator comparison later move toward a closer cross-domain articulation of that architecture. Micah is earlier on that bounded architecture delta. CFS and named precursors retain their earlier operator, causal-action, correlation-spacetime, and continuum-limit priority. Different domains and added machinery do not erase the partial match; they define its boundary.

The defensible publication claim is therefore a source-located, null-calibrated effective-equation comparison plus an atomized four-lens genealogy. Einstein-Dirac is the explicit mathematical ceiling. Temporal nonlocality and coherence are closed as bounded non-equivalences rather than left as missing work. A future gauge, scalar, kernel, or invariant comparison would require its own frozen source pair, observable, null family, and analysis plan.

\section{Appendix A. Model-Identity Chronology and Consistent-Truncation Audit for the November 2026 SIT Action (November 2026)}\label{appendix-a.-model-identity-chronology-and-consistent-truncation-audit-for-the-november-2026-sit-action-november-2026}

\subsection{A.1 Frozen July Version 80 Draft 13 versus November 2026 authoritative action}\label{a.1-frozen-july-version-80-draft-13-versus-november-2026-authoritative-action}

Sections 1--13 above were audited against the frozen \textbf{July 26, 2026 Version 80 Draft 13} (\nolinkurl{M_JUL2026_CONFORMAL}, released publicly on July 29, 2026 as Paper 002, DOI 10.5281/zenodo.21693867), which couples matter through a conformal physical metric \(\tilde{g}_{\mu\nu} = A^2(\rho_t, C) g_{\mu\nu}\). That July comparison must \textbf{not} be relabeled as a comparison against the November 2026 authoritative continuum action (\nolinkurl{M_NOV2026_AUTH_EFT_V1}, canon \texttt{v1\_1}, in the separate SIT Version 81 monograph series), which couples matter through the multiplicative portal \(+\sqrt{-g}\,f_{\mathrm{SM}}(R_{\mathrm{coh}})\mathcal{L}_{\mathrm{matter}}\) (\(f_{\mathrm{SM}}(R_{\mathrm{coh}}) = 1 + R_{\mathrm{coh}}^2/(2 M_{\mathrm{sym}}^2) \ge 1\), \(R_{\mathrm{coh}} = |\psi|/\psi_0\)) with the canonical linked-square scalar potential:

\begin{equation}
V_{\mathrm{tot}}(\rho_t, |\psi|)
= \frac{\lambda_\rho}{4}\bigl(\rho_t^2 - \rho_0^2\bigr)^2
+ \frac{\lambda_\psi}{4}\bigl(|\psi|^2 - v_0^2\bigr)^2
+ \frac{1}{2}\Bigl[m_{\mathrm{link}}(\rho_t - \rho_0) - \sqrt{\lambda_{\mathrm{portal}}}\,\bigl(|\psi|^2 - v_0^2\bigr)\Bigr]^2.
\end{equation}

\subsection{A.2 Exact consistent-truncation check for free and interacting Dirac sectors}\label{a.2-exact-consistent-truncation-check-for-free-and-interacting-dirac-sectors}

We test explicitly whether \nolinkurl{M_NOV2026_AUTH_EFT_V1} with the exact canonical potential \(V_{\mathrm{tot}}\) above admits a constant-background Einstein--Dirac sector at \((\rho_t, |\psi|, \theta, A_\mu) = (\rho_0, v_0, \theta_0, 0)\) (verified symbolically against the exact potential, Hessian, and interacting Dirac equations of motion in \nolinkurl{verify_paper048_v2.py} / \nolinkurl{paper048_v2_verification_results.json}):

\begin{enumerate}
\def\labelenumi{\arabic{enumi}.}
\tightlist
\item
  \textbf{Canonical vacuum potential, \(\mathbb{Z}_2\) lift, and mixed Hessian:} Direct evaluation of \(V_{\mathrm{tot}}\) confirms \(V_{\mathrm{tot}}(\rho_0, v_0) = 0\), \(\partial_{\rho_t} V_{\mathrm{tot}}(\rho_0, v_0) = 0\), and \(\partial_{|\psi|} V_{\mathrm{tot}}(\rho_0, v_0) = 0\) (zero classical vacuum cosmological constant), while lifting \((-\rho_0, v_0)\) by \(2 m_{\mathrm{link}}^2 \rho_0^2 > 0\). Its vacuum Hessian has nonzero portal cross-coupling \(H_{12} = -2 m_{\mathrm{link}}\sqrt{\lambda_{\mathrm{portal}}}\,v_0 \neq 0\), \(H_{11} = 2\lambda_\rho\rho_0^2 + m_{\mathrm{link}}^2\), \(H_{22} = 2(\lambda_\psi + 2\lambda_{\mathrm{portal}})v_0^2\), and strictly positive determinant \(\det H = 4\lambda_\rho\rho_0^2(\lambda_\psi + 2\lambda_{\mathrm{portal}})v_0^2 + 2 m_{\mathrm{link}}^2\lambda_\psi v_0^2 > 0\).
\item
  \textbf{Free Dirac kinetic and mass normalization:} Setting \(f_0 := f_{\mathrm{SM}}(v_0/\psi_0) \ge 1\), the free Dirac sector becomes \(\sqrt{-g}\,f_0\,\bar{\Psi}(i\gamma^\mu\nabla_\mu - m_D)\Psi\). The constant field rescaling \(\tilde{\Psi} := \sqrt{f_0}\,\Psi\) restores canonical kinetic normalization (\(1\)) identically while leaving the Dirac mass \(m_D\) strictly invariant.
\item
  \textbf{Free on-shell scalar consistency versus interacting Dirac equation of motion:} Varying the action with respect to \(\rho_t\) gives zero residual at \((\rho_0, v_0)\). Varying with respect to \(|\psi|\) yields \begin{equation}
\left.\frac{\delta S}{\delta |\psi|}\right|_{(\rho_0, v_0)} = f'_{\mathrm{SM}}(v_0)\,\mathcal{L}_m = \frac{v_0}{\psi_0^2 M_{\mathrm{sym}}^2}\,\mathcal{L}_m.
\end{equation} For a \textbf{free} minimally coupled Dirac spinor satisfying \((i\gamma^\mu\nabla_\mu - m_D)\Psi = 0\), \(\mathcal{L}_{\mathrm{Dirac}}\big|_{\mathrm{on\text{-}shell}} = 0\), so both scalar Euler--Lagrange residuals vanish identically on-shell. When an interaction is added, \textbf{the Dirac equation of motion itself changes} and one cannot hold \(\mathcal{L}_{\mathrm{Dirac}} = 0\). Specifically, for a scalar self-interaction \(\mathcal{L}_m = \mathcal{L}_{\mathrm{Dirac}} + \mathcal{L}_{\mathrm{int}}\) with \(\mathcal{L}_{\mathrm{int}} = g B^n\) and \(B := \bar{\Psi}\Psi\), variation with respect to \(\bar{\Psi}\) gives \((i\gamma^\mu\nabla_\mu - m_D)\Psi + n g B^{n-1}\Psi = 0\); contracting on the left with \(\bar{\Psi}\) yields \(\mathcal{L}_{\mathrm{Dirac}}\big|_{\mathrm{on\text{-}shell}} = -n\mathcal{L}_{\mathrm{int}}\), and therefore total on-shell matter Lagrangian and scalar source (\nolinkurl{verify_paper048_v2.py}): \begin{equation}
\begin{aligned}
& \mathcal{L}_m\big|_{\mathrm{on\text{-}shell}} = (1 - n)\mathcal{L}_{\mathrm{int}}, \\
& \left.\frac{\delta S}{\delta |\psi|}\right|_{(\rho_0, v_0),\,\mathrm{on\text{-}shell}} = (1 - n)\,f'_{\mathrm{SM}}(v_0)\,\mathcal{L}_{\mathrm{int}}.
\end{aligned}
\end{equation} Two sharp consequences follow immediately: (a) a bilinear interaction (\(n = 1\), i.e.~a mass shift \(g\bar{\Psi}\Psi\)) gives \(\mathcal{L}_{\mathrm{Dirac}} = -\mathcal{L}_{\mathrm{int}}\) and total \(\mathcal{L}_m = (1 - 1)\mathcal{L}_{\mathrm{int}} = 0\), so the constant-background truncation \textbf{still holds} (refuting any overbroad statement that every nonzero interaction breaks it); and (b) for a four-fermion interaction (\(n = 2\), \(\mathcal{L}_{\mathrm{int}} = \lambda B^2\)), \(\mathcal{L}_{\mathrm{Dirac}} = -2\lambda B^2\), \(\mathcal{L}_m = -\lambda B^2\), and the scalar source is \textbf{\(-f'_{\mathrm{SM}}(v_0)\,\mathcal{L}_{\mathrm{int}} = -\frac{v_0}{\psi_0^2 M_{\mathrm{sym}}^2}\lambda B^2\)} (opposite in sign to the naive \(+f'_0\mathcal{L}_{\mathrm{int}}\) obtained by erroneously keeping \(\mathcal{L}_{\mathrm{Dirac}} = 0\)).
\item
  \textbf{Einstein stress-energy and gauge-current normalization:} Varying \(f_0\sqrt{-g}\mathcal{L}_{\mathrm{Dirac}}(\Psi)\) with respect to the tetrad \emph{prior} to applying the on-shell identity yields \(f_0 T^{\mathrm{Dirac}}_{\mu\nu}(\Psi) = T^{\mathrm{Dirac}}_{\mu\nu}(\tilde{\Psi})\), while scalar kinetic and potential contributions vanish identically at constant \((\rho_0, v_0)\). In the vector sector, with the unrescaled fermion \(\Psi\) inside \(\sqrt{-g}\,f_0\mathcal{L}_m(\Psi)\), varying with respect to \(A_\mu\) produces the dressed gauge-field source \(f_0 q_\Psi \bar{\Psi}\gamma^\mu\Psi\) (before any additional gauge-field normalization \(f_{\mathrm{EM}}(R_0)\)); calling \(J_{\mathrm{gauge}}^\mu = q_\Psi \bar{\Psi}\gamma^\mu\Psi\) the bare gauge current leaves \(f_0\) explicit in the field equation, whereas canonical fermion normalization \(\tilde{\Psi} = \sqrt{f_0}\,\Psi\) absorbs \(f_0\) identically into \(q_\Psi\bar{\tilde{\Psi}}\gamma^\mu\tilde{\Psi}\) (\nolinkurl{verify_paper048_v2.py}). In both conventions, the fermion-number current \(J_F^\mu = \bar{\Psi}\gamma^\mu\Psi\) is distinct from the charge-weighted gauge source: a \textbf{neutral} Dirac field (\(q_\Psi = 0\)) can carry nonzero \(J_F^\mu\) without sourcing \(A_\mu\), whereas a charged Dirac field (\(q_\Psi \neq 0\)) sources \(A_\mu\) and breaks the \(A_\mu = 0\) background.
\end{enumerate}

\textbf{Preserved zero-discrimination boundary:} This calculation establishes continuity of the shared constant-background free (and \(n=1\) bilinear) Einstein--Dirac effective limit from July Version 80 to November \nolinkurl{M_NOV2026_AUTH_EFT_V1}. Exactly as in Section 10, that shared limit carries \textbf{zero discriminatory value} between Causal Fermion Systems and SIT, and does not derive either framework from the other.

\section{Status Note and Version Lineage}\label{status-note-and-version-lineage}

The July 28, 2026 edition was publicly deposited on July 29, 2026 as Paper 048 (Zenodo DOI 10.5281/zenodo.21693926, \url{https://pub.svgn.io/papers/paper-48/}). This November 20, 2026 Version 2 candidate preserves the July four-lens audit and frozen July 26, 2026 Version 80 Draft 13 baseline intact while adding Appendix A to separate the July and November SIT actions and execute the exact canonical-potential and interacting-Dirac truncation audit (\nolinkurl{verify_paper048_v2.py}). It remains a private review candidate awaiting explicit author release authorization.

\section{References}\label{references}

\begin{enumerate}
\def\labelenumi{\arabic{enumi}.}
\tightlist
\item
  Finster, F., Kindermann, S., and Treude, J.-H. (2025). \emph{Causal Fermion Systems: An Introduction to Fundamental Structures, Methods and Applications}. arXiv:2411.06450v2.
\item
  Finster, F., Kleiner, J., and Paganini, C. F. (2024). Causal Fermion Systems as an Effective Collapse Theory. \emph{Journal of Physics A: Mathematical and Theoretical}, 57, 395303. \url{doi:10.1088/1751-8121/ad6d24}.
\item
  Fischer, P., and Paganini, C. F. (2025). \emph{Causal Fermion Systems: Spacetime as the Web of Correlations of a Many-Body Quantum System}. arXiv:2504.19272v1. \url{https://arxiv.org/abs/2504.19272}.
\item
  Farnsworth, S., Finster, F., Paganini, C. F., and Singh, T. P. (2026). \emph{Causal Fermion Systems, Non-Commutative Geometry and Generalized Trace Dynamics}. arXiv:2603.05018v2. \url{https://arxiv.org/abs/2603.05018}.
\item
  Fischer, P., and Finster, F. (2026). The Continuum Limit Analysis of Causal Fermion Systems for Curved Spacetimes. arXiv:2605.30199v1.
\item
  Blumberg, M. (2026). \emph{Super Information Theory}, Version 80 Draft 13 current manuscript.
\item
  Blumberg, M. (2025). \emph{A Constructive Analysis of Convergence: Causal Fermion Systems and Super Information Theory}. Preserved private source dossier.
\item
  Blumberg, M. (2017). \emph{The Neural Lace Podcast Origin: The Neo Mind Cycle Concept}. Silicon Valley Global News, published 11 April 2017. \url{https://www.youtube.com/watch?v=YrX_68oKuVs}.
\item
  Blumberg, M., and Sanders, B. (2017). \emph{The Neural Lace Podcast \#2}. Silicon Valley Global News, published 13 April 2017. \url{https://www.youtube.com/watch?v=UpOzHjKGa0g}.
\item
  Blumberg, M., and Matzner, E. (2017). \emph{The Neural Lace Podcast \#3}. Silicon Valley Global News, published 22 April 2017. \url{https://www.youtube.com/watch?v=_yKjtTVoVlU}.
\item
  Blumberg, M. (2022). \texttt{a0010z.md}, public GitHub commit \nolinkurl{40c0b37fc9376a6630ad9d3b09d911997c1c17f4}, 10 June 2022. \url{https://github.com/v5ma/selfawarenetworks/commit/40c0b37fc9376a6630ad9d3b09d911997c1c17f4}.
\item
  Blumberg, M. (2022). \texttt{a0387z.md}, public GitHub commit \nolinkurl{4c451f016cd076bcb635f4305a55efe54bbf4870}, 12 June 2022. \url{https://github.com/v5ma/selfawarenetworks/commit/4c451f016cd076bcb635f4305a55efe54bbf4870}.
\item
  Blumberg, M. (2022). \texttt{a0322z.md}, public GitHub commit \nolinkurl{4e0cf1ee0975c01dc57182368497d3f0cd864030}, 28 November 2022. \url{https://github.com/v5ma/selfawarenetworks/commit/4e0cf1ee0975c01dc57182368497d3f0cd864030}.
\end{enumerate}

\end{document}
